On considère une solution aqueuse $(\mathrm{S_{0}})$ d'ammoniac $\ce{NH3}$, de
volume $V_{0}$ et de concentration molaire $C_{0}=\qty{1.0e-2}{\mol\per\L}$. Le
$\mathrm{pH}$ de cette solution à $\qty{25}{\degreeCelsius}$ vaut
$\mathrm{pH}=10,6$.

L'équation de la réaction modélisant la transformation entre l'ammoniac et l'eau
est~:
\[
    \ce{NH_{3(aq)} + H2O_{(\ell)} <=> NH^{+}_{4(aq)} + HO^{-}_{(aq)}}
\]

\begin{donnees}
    \begin{itemize}
        \item Produit ionique de l'eau à $\qty{25}{\degreeCelsius}$~:
              $\mathrm{K_{e}}=10^{-14}$.
    \end{itemize}
\end{donnees}
\begin{enumerate}
    \item Montrer que la concentration molaire effective des ions ammonium
          $\ce{NH^+_{4(aq)}}$à l'état d'équilibre du système est donnée par la
          relation:
          $\bigl[\ce{NH^+_{4(aq)}}\bigr]_{\eq}=
          \frac{\mathrm{K_{e}}}{10^{-\mathrm{pH}}}$.
          Calculer sa valeur.

    \item Calculer la valeur du quotient de réaction $Q_{r,\eq}$ du système
          chimique à l'équilibre. En déduire la valeur de la constante
          d'équilibre $\mathrm{K}$ associée à l'équation de la réaction.

    \item La constante d'acidité du couple $(\ce{NH^+_{4(aq)}}/\ce{NH_{3(aq)}})$
          s'exprime par~: $\mathrm{K_A}=\frac{\mathrm{K_e}}{\mathrm{K}}$.

          Calculer le $\mathrm{pK_A}$ de ce couple.

    \item On mélange un volume de la solution $(\mathrm{S_{0}})$ d'ammoniac avec
          un volume d'une solution de chlorure d'ammonium
          $\ce{NH^+_{4(aq)} + Cl^-_{(aq)}}$. Le $\mathrm{pH}$ du mélange est
          $\mathrm{pH}=6,2$.
    
          Tracer le diagramme de prédominance des espèces du couple
          $(\ce{NH^+_{4(aq)}}/\ce{NH_{3(aq)}})$. En déduire l'espèce
          prédominante de ce couple dans le mélange.
\end{enumerate}

